\chapter{Introduction to Web Development with Python (Flask)}

\section{Introduction: Real-Life Hook}

Think about things students see online every day:

\begin{itemize}
    \item School websites with notices  
    \item Online quizzes or forms  
    \item Simple blogs or portfolio websites  
\end{itemize}

Flask is a \textbf{Python framework that lets you create your own web applications} easily and quickly.

\section{Concept: What is Flask?}

\begin{definition}{Flask}
Flask is a \textbf{lightweight Python web framework} that allows you to build web applications with minimal code.
\end{definition}

- Lightweight means it is \textbf{easy for beginners}  
- Provides \textbf{routes, templates, and server handling}  
- Perfect for \textbf{small projects, prototypes, and learning web development}

\section{Setting Up Flask}

Install Flask:

\begin{verbatim}
pip install Flask
\end{verbatim}

\section{Creating Your First Flask App}

\begin{verbatim}
from flask import Flask

app = Flask(__name__)

@app.route("/")
def home():
    return "Hello, Flask!"

if __name__ == "__main__":
    app.run(debug=True)
\end{verbatim}

Explanation:

\begin{itemize}
    \item \texttt{Flask(\_\_name\_\_)} creates the app  
    \item \texttt{@app.route("/")} defines the home page  
    \item \texttt{app.run(debug=True)} runs the web server  
\end{itemize}

\section{Using HTML Templates}

Flask can render HTML for nicer web pages:

1. Create a folder \texttt{templates/}  
2. Add \texttt{index.html}:

\begin{verbatim}
<!DOCTYPE html>
<html>
<head><title>Home Page</title></head>
<body>
<h1>Welcome to My Website!</h1>
<p>This is a Flask web page.</p>
</body>
</html>
\end{verbatim}

3. Update \texttt{app.py}:

\begin{verbatim}
from flask import Flask, render_template

app = Flask(__name__)

@app.route("/")
def home():
    return render_template("index.html")

if __name__ == "__main__":
    app.run(debug=True)
\end{verbatim}

- \texttt{render\_template()} loads HTML files from \texttt{templates/} folder  

\section{Handling Forms and User Input}

\begin{verbatim}
from flask import Flask, render_template, request

app = Flask(__name__)

@app.route("/", methods=["GET", "POST"])
def home():
    message = ""
    if request.method == "POST":
        name = request.form["name"]
        message = f"Hello, {name}!"
    return render_template("form.html", message=message)

if __name__ == "__main__":
    app.run(debug=True)
\end{verbatim}

\texttt{form.html}:

\begin{verbatim}
<!DOCTYPE html>
<html>
<body>
<h1>Enter Your Name</h1>
<form method="POST">
  <input type="text" name="name" placeholder="Your Name">
  <input type="submit" value="Submit">
</form>
<p>{{ message }}</p>
</body>
</html>
\end{verbatim}

- Students learn \textbf{GET/POST methods, input handling, and dynamic web pages}  

\section{Mini-Project: Student Feedback Web App}

Concept:

\begin{itemize}
    \item Students enter their name and feedback in a form  
    \item Server stores or displays the feedback  
    \item Page shows a thank-you message after submission  
\end{itemize}

Basic structure:

\begin{verbatim}
from flask import Flask, render_template, request

app = Flask(__name__)
feedbacks = []

@app.route("/", methods=["GET", "POST"])
def home():
    if request.method == "POST":
        name = request.form["name"]
        feedback = request.form["feedback"]
        feedbacks.append({"name": name, "feedback": feedback})
        return render_template("thankyou.html", name=name)
    return render_template("feedback_form.html")

if __name__ == "__main__":
    app.run(debug=True)
\end{verbatim}

- Students learn \textbf{dynamic content, lists/dictionaries to store data, and templates}  

\section{Step-by-Step Diagram: Flask App Workflow}

\begin{center}
\begin{tikzpicture}[node distance=1.5cm]
\node (start) [draw, ellipse, fill=green!20] {User Opens Web Page};
\node (request) [draw, rectangle, fill=yellow!20, below of=start] {HTTP Request Sent to Flask Server};
\node (route) [draw, rectangle, fill=orange!20, below of=request] {Server Processes Route / Logic};
\node (template) [draw, rectangle, fill=yellow!20, below of=route] {Render HTML Template};
\node (response) [draw, rectangle, fill=orange!20, below of=template] {Send Response Back to User};
\node (end) [draw, ellipse, fill=green!20, below of=response] {Page Displayed in Browser};
\draw[->] (start) -- (request);
\draw[->] (request) -- (route);
\draw[->] (route) -- (template);
\draw[->] (template) -- (response);
\draw[->] (response) -- (end);
\end{tikzpicture}
\end{center}

\section{Key Concepts}

\begin{itemize}
    \item Flask is a lightweight Python framework for \textbf{web development}  
    \item Routes map \textbf{URLs to functions}  
    \item Templates allow \textbf{dynamic HTML content}  
    \item Forms and POST requests allow \textbf{user input handling}  
    \item Students can create \textbf{interactive web apps} with minimal code  
\end{itemize}

\section{End-of-Chapter Exercises}

\begin{exercise}{Very Short Questions}
\begin{enumerate}
    \item What is Flask used for?  
    \item Name two types of HTTP requests.  
    \item What folder is used to store HTML templates in Flask?  
\end{enumerate}
\end{exercise}

\begin{exercise}{Short Answer Questions}
\begin{enumerate}
    \item Write a Flask route that returns “Hello World”.  
    \item Explain the difference between GET and POST methods.  
    \item How do you render an HTML file in Flask?  
\end{enumerate}
\end{exercise}

\begin{exercise}{Activity-Based Questions}
\begin{enumerate}
    \item Create a Flask app that accepts student names and marks, then displays the highest mark.  
    \item Build a small quiz web app with 3 questions and display the score after submission.  
    \item Modify the feedback form to show all previous feedbacks below the form.  
\end{enumerate}
\end{exercise}

\section{Teacher Notes}

\marginnote{
\textbf{Learning Objectives:}  
\begin{itemize}
    \item Students understand Flask basics and routes  
    \item Students can create HTML templates and forms  
    \item Students can build simple interactive web applications  
\end{itemize}

\textbf{Teaching Tips:}  
\begin{itemize}
    \item Start with simple “Hello World” route before forms  
    \item Explain HTTP requests and responses in simple terms  
    \item Encourage students to customize HTML with colors and fonts  
\end{itemize}

\textbf{Assessment:}  
\begin{itemize}
    \item Check correct route handling and template rendering  
    \item Observe user input handling and dynamic content display  
\end{itemize}
}